\textbf{结论名称：} 倒序相加法（首末等距项和为定值）

\textbf{适用条件：} 数列 $\{a_n\}$ 中与首末两项等距离的两项之和为定值，即存在常数 $C$ 使
\[
a_k + a_{n+1-k} = C \quad (k=1,2,\dots,n)
\]
成立。最典型的特例是等差数列，此时 $C = a_1 + a_n$。

\textbf{变量说明：} $a_n$ 为数列通项，$S_n = \sum_{k=1}^n a_k$ 为前 $n$ 项和，$C$ 为等距项和的定值常数。

\textbf{核心公式：} 
倒序写出 $S_n$ 并与原序相加：
\[
2S_n = \sum_{k=1}^n (a_k + a_{n+1-k}) = nC.
\]
当 $C = a_1 + a_n$（如等差数列）时，得到常用求和公式
\[
\boxed{S_n = \frac{n(a_1 + a_n)}{2}}.
\]

\textbf{使用场景：} 等差数列的求和推导；满足 $f(x)+f(1-x)=c$ 的函数值构成的数列求和（如 $f(k/n)$ 型）；对称式数列的快速求和。

\textbf{简单例子：} 
计算 $1+2+\cdots+100$。
令 $S = 1+2+\cdots+100$，倒序写 $S = 100+99+\cdots+1$。
两式相加得 $2S = (1+100)+(2+99)+\cdots+(100+1) = 101 \times 100$，
所以 $S = \dfrac{101 \times 100}{2} = 5050$。直接验证 $\dfrac{100(1+100)}{2}=5050$。